% BHU PYQ -- Manually transcribed & error-corrected
% Paper: BCF-113 : Business Economics
% Exam: B.Com. (Hons.) FMM Semester I Examination, 2022-23

\documentclass[11pt,a4paper]{article}
\usepackage[a4paper,top=2.5cm,bottom=2.5cm,left=2.8cm,right=2.8cm,headheight=14pt]{geometry}
\usepackage{amsmath,amssymb}
\usepackage[T1]{fontenc}
\usepackage[utf8]{inputenc}
\usepackage{lmodern,microtype,fancyhdr,enumitem,hyperref,lastpage,parskip}

\hypersetup{
    colorlinks=true,
    urlcolor=black,
    linkcolor=black,
    pdftitle={BCF-113: Business Economics -- BHU B.Com. (Hons.) FMM Sem I 2022-23}
}

\pagestyle{fancy}
\fancyhf{}
\renewcommand{\headrulewidth}{0.4pt}
\renewcommand{\footrulewidth}{0.4pt}
\fancyhead[L]{\small   \textit{Banaras Hindu University}}
\fancyhead[R]{\small   \textit{BCF-113: Business Economics}}
\fancyfoot[C]{\small Page  \thepage\ of \pageref{LastPage}}


\newlist{parts}{enumerate}{2}
\setlist[parts,1]{label=(\alph*),leftmargin=*,itemsep=5pt,topsep=3pt}
\setlist[parts,2]{label=(\roman*),leftmargin=*,itemsep=3pt,topsep=2pt}

\newcommand{\pts}[1]{\hfill   \textit{[#1]}}


\begin{document}


\begin{center}
  {\large\bfseries BANARAS HINDU UNIVERSITY}\\[2pt]
  {
\normalsize B.Com. (Hons.) FMM --- Semester I Examination, 2022-23}\\[6pt]
  \rule{0.8\linewidth}{0.5pt}\\[6pt]
  {\large\bfseries Paper No. BCF-113: Business Economics}\\[4pt]
  {
\normalsize Time: 3 Hours\hfill Maximum Marks: 70}
\end{center}

\rule{\linewidth}{0.4pt}

\smallskip



\noindent   \textit{Instructions: Answer all questions. All questions carry equal marks (14 marks each).}

\bigskip




\noindent   \textbf{Question 1.}
What are the central problems in an economy? How can those problems be solved in different economic structures?\pts{5+9}

\centerline{   \textbf{OR}}

Differentiate between microeconomics and macroeconomics as a discipline. What is the role of price theory?\pts{9+5}

\medskip\rule{\linewidth}{0.2pt}




\noindent   \textbf{Question 2.}
Answer the following questions:

\begin{parts}
  \item Why is a demand curve downward sloping?\pts{7}
  \item Why is an indifference curve downward sloping?\pts{7}
\end{parts}

\centerline{   \textbf{OR}}

Explain the Revealed Preference Theory.\pts{14}

\medskip\rule{\linewidth}{0.2pt}




\noindent   \textbf{Question 3.}
Answer the following questions:

\begin{parts}
  \item ``LAC is an envelope of various SAC curves.'' Explain.\pts{9}
  \item Why does a short-run average cost curve have a U-shape?\pts{5}
\end{parts}

\centerline{   \textbf{OR}}

Answer the following questions:

\begin{parts}
  \item Discuss the modern theory of factor pricing.\pts{8}
  \item Explain the marginal productivity theory of distribution.\pts{6}
\end{parts}

\medskip\rule{\linewidth}{0.2pt}




\noindent   \textbf{Question 4.}
Discuss consumer equilibrium using Ordinal Utility Theory. Explain whether it provides the same result as Cardinal Utility Analysis.\pts{9+5}

\centerline{   \textbf{OR}}

What are the properties of an Isoquant? Derive the producer's equilibrium.\pts{5+9}

\medskip\rule{\linewidth}{0.2pt}




\noindent   \textbf{Question 5.}
Among perfect competition and monopoly market structures, which one is more efficient and why? Explain.\pts{14}

\centerline{   \textbf{OR}}

Discuss the monopolistic market structure.\pts{14}

\vfill
\rule{\linewidth}{0.4pt}

\begin{center}\small
\href{https://bhu-exam.vercel.app/}{Click here to visit our website}\\[4pt]
\href{https://me-aryan.vercel.app/}{Click here to visit our contributor}\\[4pt]
\href{https://quashan-me.vercel.app/}{Click here to visit our contributor}
\end{center}

\end{document}
